\chapter{Theory}\label{ch:theory}

This chapter explores the theory behind important concepts such as KDE, compares different smoothing methods, and examines a simple reference implementation.

\section{Kernel Density Estimation}\label{sec:kde}

\subsection{Definition}

In order to fully understand the following sections, it is crucial to understand the concept of Kernel Density Estimation (KDE), since it is used in many of the smoothing methods described below.

If ($x_1$, $x_2$, \dots, $x_n$) are samples drawn from some univariate distribution with an unknown density function $f(x)$ at any given point $x$, we can estimate the shape of the density function using a \emph{kernel density estimator} given by
\[
\hat{f}_{h}(x)=\frac{1}{n}\sum_{i=1}^n K_{h}(x-x_{i})=\frac{1}{nh}\sum_{i=1}^n K\left( \frac{x-x_{i}}{h} \right)
\]

where $K$ is the kernel, a non-negative function, and $h > 0$ is a smoothing parameter called the \emph{bandwidth}.\ \cite{kde:01}

\subsection{Kernel functions}\label{subsec:kernelfunctions}

The kernel used in a KDE governs the smoothing behavior of the estimation. As seen in the formula, the distance between the given point $x$ and each sample point $x_i$ is passed through the kernel. A kernel function is a function $K$ that is non-negative everywhere, normalized, and symmetric, which means that:
\[
\int_{-\infty}^\infty K(u)\,du=1 \text{ and } K(u)=K(-u) \text{ for all }u
\]

The parameter $u$ here is often chosen to represent the standardized distance:
\[
u=\frac{x-x_i}{h}
\]
\nolinebreak\cite{kde:book}

The Epanechnikov or parabolic kernel
\[
K(u)=\frac{3}{4}(1-u^{2}) \text{ for }|u| \leq 1,\quad K(u)=0 \text{ for } |u| > 1
\]
is the ``ideal'' estimator because it minimizes the mean squared error (MSE) for any given bandwidth\ \cite{kde:book}.

However, by far the most common kernel is $\phi$, the probability density function of a normal distribution with $\mu=0$ and $\sigma=1$\footnote{Also sometimes called the \emph{Gaussian kernel}}. Plugging in $\phi$ for $K$, the KDE simplifies to:
\[
\hat{f}_{h}(x)=\frac{1}{nh\sigma} \frac{1}{\sqrt{ 2\pi }}\sum_{i=1}^n \exp\left( \frac{-{(x-x_{i})}^{2}}{2h^{2}\sigma^{2}} \right)
\]

Many formulas also use $K_h$, which is the kernel function with the bandwidth already applied.
\[
K_h(x)=\frac{1}{h}K\left(\frac{x}{h}\right)
\]


\subsection{Example}

\begin{figure}[H]
	\centering
    \includegraphics[scale=0.35]{images/kdeplot.png}
	\caption{Kernel density estimations with different bandwidths}\label{fig:KDEPlot}
\end{figure}

\autoref{fig:KDEPlot} shows the original normal distribution and three kernel density estimations with different bandwidths (generated using a Python script).
\begin{itemize}
    \item Red: $h=0.4$
    \item Green: $h=1$
    \item Blue: $h=0.15$
    \item Gray: the original normal distribution
\end{itemize}
and using $\phi$ as the kernel.

\subsection{Bandwidth selection}

As seen in \autoref{fig:KDEPlot}, the bandwidth has a great impact on the final output of the estimation. Calculating the ideal bandwidth for any sample is not possible because doing so would require the original density function $f$ or its second derivative $f''$ \ \cite{kde:book}. There are, however, approximations, a few of which are covered here.

\subsubsection{Scott method}
The Scott method or Scott's Rule works simply by estimating
\[
h \approx n^{-1/(d+4)}
\]
where $n$ is the number of data points and $d$ is the number of dimensions. For one-dimensional data, this becomes
\[
h = n^{-1/5}
\]
\ \cite{kde:03}

\subsubsection{Silverman method}

Provided that a Gaussian basis function is used as the kernel and the underlying density is also Gaussian, the Silverman method is ideal (that is, it minimizes the mean squared error).
\[
h={\left(\frac{4\hat{\sigma}^5}{3n}\right)}^{1/5}\approx 1.06\hat{\sigma}n^{-1/5}
\]
where $\hat{\sigma}$ is the standard deviation of the sample.

This can be further improved by replacing $\hat{\sigma}$ with
\[
A=\min\left(\hat{\sigma}, \frac{IQR}{1.34}\right)
\]
where $IQR$ is the interquartile range. Another modification is to reduce the scaling factor from $1.06$ to $0.9$. The final formula then is
\[
h=0.9\,\min\left(\hat{\sigma}, \frac{IQR}{1.34}\right)\,n^{-1/5}
\]
and is also known as \emph{Silverman's rule of thumb}.\ \cite{kde:book2}

\subsection{Weighted KDE}

The weighted KDE extends the KDE with the ability to associate a weight with each data point. The general formula for weighted KDE is
\[
\hat{f}(x)=\frac{1}{\sum_i w_i}\sum_{i=1}^n w_{i} K_{h}(x-X_{i})
\]
adapted from\ \cite{kde:weighted01} where $w_i$ is the associated weight of each data point.

When using a method to calculate the bandwidth $h$ using a weighted KDE, the number of samples must be replaced with $n_\mathrm{eff}$, the effective sample size, estimated by
\[
n_\mathrm{eff}=\frac{{\left(\sum_{i=1}^n w_i\right)}^2}{\sum_{i=1}^n w_i^2}
\]
\nolinebreak\cite{kde:book3}

\subsection{Spatiotemporal KDE}

The Spatiotemporal KDE adapts the KDE to use location and time, so instead of $\hat{f}_h(x)$ it is:
\[
\hat{f}(x,y,t)=\frac{1}{n_\mathrm{eff}\cdot h_{x}\cdot h_{y}\cdot h_{t}}\sum_{i=1}^n w_{i}\cdot K\left( \frac{x-x_{i}}{h_{x}} \right)\cdot K\left( \frac{y-y_{i}}{h_{y}} \right)\cdot K\left( \frac{t-t_{i}}{h_{t}} \right)
\]
where
\begin{itemize}
    \item $\hat{f}(x,y,t)$ is the estimated probability density for a location in space and time
    \item $n_\mathrm{eff}$ is the effective number of points, calculated using Kish's effective sample size (see\ \cite{kde:book3})
    \item $h_x$, $h_y$ and $h_t$ are the bandwidths in each dimension
    \item $x_i$, $y_i$ and $t_i$ are the coordinates of each data point, $w_i$ is the optional weight
    \item $K$ is a kernel function (see \autoref{subsec:kernelfunctions})
\end{itemize}


\section{Different Smoothing Methods}\label{sec:smoothingmethods}
This section will explain and evaluate different smoothing methods.

In general, a spatiotemporal smoothing method provides a function, often called $\hat{z}(x_0)$, that returns an interpolated value at a specific location $x_0$.

Many methods can be represented using a weighted sum using the general formula
\[
\hat{z}(x_0)=\sum_{i=1}^N \lambda_i \, z(x_i)
\]

where
\begin{itemize}
    \item $x_i$ is a measured data point.
    \item $z(x_i)$ is the value at the measured point.
    \item $\lambda_i$ is a weight associated with the data point. Depending on the method, $\sum_{i=1}^N \lambda_i$ may equal one, making the method a weighted average.
    \item $N$ is the total number of data points.
\end{itemize}

Unless specified otherwise, the following smoothing methods will use this formula for interpolation, with the specific method usually being defined by the choice of $\lambda_i$.

Oftentimes, $\lambda_i$ is also a function of the distance $d_{s_i}$ between the interpolation point $x_0$ and a data point $x_i$. For the following methods, the Euclidean distance is used, which means that
\[
d_{s_i}=\sqrt{{(x_i-x_0)}^2+{(y_i-y_0)}^2}
\]

Using the definitions provided for $\hat{z}(x_0)$ and $d_{s_i}$ allows for the definition of a smoothing method by providing $\lambda_i$ or $\hat{z}(x_0)$ if necessary.

\subsection{Nadaraya-Watson kernel regression}\label{subsec:nw_regression}

The NW\footnote{Nadaraya-Watson} kernel regression, or NW estimator, was first proposed by Nadaraya and Watson in 1964 (see\ \cite{nw:nadaraya} and\ \cite{nw:watson}). It can be seen as a weighted average in which each data point is weighted based on its distance to the interpolation point.

The derivation (see\ \cite{nw:book} and\ \cite{nw:web01}) assumes that data points are of the form $\{(X_1,Y_1), \dots,(X_n,Y_n)\}$. The goal is to find the \emph{conditional expectation} $E$ of $Y$ relative to $X$, which takes the form
\[
E(Y|X=x_0)=\int \frac{yf(y,x)}{h(x)}\,dy
\]
where $f(y,x)$ is the joint density function of $X$ and $Y$ and  $h(x)$ is the \emph{marginal density} (see\ \cite{nw:mdense}) of $f(y,x)$, so
\[
h(x)=\int f(y,x)\,dy=f_X(x)
\]

$h(x)$ is then replaced with a KDE (see \autoref{sec:kde})
\[
\hat{h}(x)=\frac{1}{n}\sum_{j=1}^n K_h(x-X_j)
\]
and so is $f(y,x)$
\[
\hat{f}(y,x)=\frac{1}{n}\sum_{i=1}^n K_h(x-X_i) K_h(y-Y_i)
\]
which are substituted into the formula and simplified:
\begin{align*}
    \hat{E}(Y|X) &= \int \frac{y\hat{f}(y,x)}{\hat{h}(x)}\,dy\\
    &=\int \frac{y \frac{1}{n}\sum_{i=1}^n K_h\left(x-X_i\right) K_h\left(y-Y_i\right)}{\frac{1}{n}\sum_{j=1}^n K_h\left(x-X_j\right)}\,dy\\
    &=\frac{\frac{1}{n}\sum_{i=1}^n K_h\left(x-X_i\right)\int y K_h\left(y-Y_i\right)\,dy}{\frac{1}{n}\sum_{j=1}^n K_h\left(x-X_j\right)}\\
    &=\frac{\frac{1}{n}\sum_{i=1}^n K_h\left(x-X_i\right)Y_i}{\frac{1}{n}\sum_{j=1}^n K_h\left(x-X_j\right)}\\
    &=\sum_{i=1}^n \frac{K_h\left(x-X_i\right)Y_i}{\sum_{j=1}^n K_h\left(x-X_j\right)}
\end{align*}
\vspace{1.5em}

This is the NW estimator. To fit the format of this thesis, substitute $Y$ with $z(x)$ and $x - X_i$ with $d_{s_i}$
\[
\hat{z}(x_0)=\sum_{i=1}^N \frac{K_h\left(d_{s_i}\right)\cdot z\left(x_i\right)}{\sum_{j=1}^N K_h\left(d_{s_j}\right)}
\]

To fit the general formula from \autoref{sec:smoothingmethods}, the weighting factor $\lambda_i$ is extracted from the estimator:
\[
\lambda_i(d_{s_i})=\frac{K_h\left(d_{s_i}\right)}{\sum_{j=1}^N K_h\left(d_{s_j}\right)}
\]

\subsection{Inverse Distance Weighting}\label{subsec:idw}

IDW\footnote{Inverse Distance Weighting} was first proposed by Shepard in 1968 (see\ \cite{idw:paper01}). Classic IDW can be described as a simple weighted sum over all data points, using the inverse of the distance to the interpolation point. It is equivalent to the setup described in \autoref{sec:smoothingmethods} with
\[
\lambda_i(d_{s_i})=\frac{d_{s_i}^{-u}}{\sum_{j=1}^N d_{s_j}^{-u}}
\]  
where $u$ is a parameter that indicates how strongly the weight decreases as the distance increases.

\emph{Improved} IDW only considers points within a certain radius $R$ of the interpolation point, greatly increasing performance with little loss in quality because distant points have almost no influence. Performance can be further improved if a spatial data structure such as a quadtree or k-d tree is used. Although the parameter $R$ has a major impact on the final result, no empirical approach has been developed to determine its optimal value\ \cite{mdpi:01}.

\begin{figure}[H]
	\centering
    \includegraphics[scale=0.75]{images/R.pdf}
	\caption{Improved IDW\@. Only the points in blue are included in the interpolation\\(see\ \cite{mdpi:01})}\label{fig:IDW}
\end{figure}

There can also be cases where $d_{s_i}$ is zero for a data point, which means that it coincides with the interpolation point. The standard solution is for $\hat{z}$ to return the value of that data point or, if necessary, the mean of all data points equal to the interpolation point.\footnote{Taking the average of all matching data points is not addressed in the sources but is included here because our data can contain duplicate points with different values}

\subsection{Distribution-Based Distance Weighting}\label{subsec:ddw}

DDW\footnote{Distribution-Based Distance Weighting} (see\ \cite{mdpi:01}) is an extension of the Improved IDW approach, which replaces purely distance-based weighting with a weight function that includes the distribution of point values around the interpolation point. This allows the method to adapt to different environments, potentially providing smoother and more context-sensitive interpolation behavior compared to IDW\@.

This is achieved by using the Improved IDW setup and replacing $\lambda_i$ with
\[
\lambda_i(d_{s_i}) = \frac{D(d_{s_i}, \mu_R, k\sigma_R)}{\sum_{j=1}^N D(d_{s_j}, \mu_R, k\sigma_R)},
\]
where
\begin{itemize}
    \item $D$ is a chosen probability distribution
    \item $R$ is the radius of a circle centered at the interpolation point. Only data points within this circle are considered for interpolation (see \autoref{fig:IDW})
    \item $\sigma_R$ is the standard deviation of all data values scaled by the spread
    \item $\mu_R$ is always 0
    \item $k$ is a user-defined scaling factor applied to the standard deviation\footnote{The parameter $k$ was introduced in this thesis.}
\end{itemize}

Note that the original paper does not include the sum over all weights in the denominator, but it was added in this thesis because it is required for the weights to add up to one.

The exact formula for $\sigma_R$ is
\[
\sigma_R=R\cdot\frac{\sigma}{\max(|\min(z)-\mu_z|,|\max(z)-\mu_z|)}
\]
where
\begin{itemize}
    \item $\sigma$ is the standard deviation of all point values within the radius $R$
    \item $z$ is the list of all values within the radius $R$
    \item $\mu_z$ is the mean of all values within the radius $R$
\end{itemize}

This has the effect of making the distance more important when the values are similar and increasing the influence of more distant points when the values have a higher variance.

This thesis examines three candidate distributions: the Normal, Cauchy, and Laplace distributions. These were selected because they
\begin{itemize}
    \item are symmetric around their location parameter;
    \item are continuous;
    \item are not confined to a finite interval;
    \item feature a parameter that controls the location (denoted here as $\mu$);
    \item include a parameter that influences scale or spread;
    \item exhibit distinct shapes, which makes it possible to evaluate DDW with different distributional profiles.
\end{itemize}

The probability density functions used are
\[
f(x,\sigma)=\frac{1}{\sqrt{2\pi\sigma^2}}e^{-\frac{{(x-\mu)}^2}{2\sigma^2}}
\]
for the Normal distribution,
\[
f(x,\Gamma)=\frac{\Gamma}{2\pi}\cdot\frac{1}{{(x-\mu)}^2+{\left(\frac{\Gamma}{2}\right)}^2}
\]
for the Cauchy distribution, and
\[
f(x,b)=\frac{1}{2b}e^{-\frac{|x-\mu|}{b}}
\]
for the Laplace distribution.

Note that each distribution uses a different convention for its scale parameter. Only the scale parameter of the Normal distribution corresponds directly to the standard deviation. In order to maintain comparability between distributions and to provide the user with additional control over effective spread, the parameter $k$ is introduced to rescale the standard deviation. Therefore, a complete specification of $\lambda_i$ now requires both the selection of a distribution and a value for $k$.\ \cite{mdpi:01}

\subsection{Kriging}\label{subsec:kriging}

Kriging was developed in 1951 by Danie G. Krige to estimate the spatial distribution of gold at the Witwatersrand reef in South Africa. The theoretical basis was formalized by Georges Matheron in 1960 in his paper ``Krigeage d'un panneau rectangulaire par sa périphérie''.\ \cite{kriging:invention, kriging:invention2}

According to\ \cite{shape:book01}, there are more than twenty different kriging methods. The two most commonly used methods are Ordinary kriging and Universal kriging. The ArcGIS Pro documentation describes the differences between them:

\begin{quote}
    \textit{Ordinary kriging is the most general and widely used of the kriging methods and is the default. It assumes the constant mean is unknown. This is a reasonable assumption unless there is a scientific reason to reject it.\\Universal kriging assumes that there is an overriding trend in the data [\ldots] and it can be modeled by a deterministic function. [\ldots] Universal kriging should only be used when you know there is a trend in your data and you can give a scientific justification to describe it.}\ \cite{kriging:web01}
\end{quote}

In this thesis, only Ordinary kriging is used, since the available data do not exhibit an identifiable deterministic trend. Other kriging types are not discussed here to maintain a focused scope for this chapter.

The first step in kriging is to find the spatial dependence (spatial autocorrelation) of the data. This is achieved by creating a semivariogram and fitting a mathematical model to it.

A semivariogram consists of all pairs of data points and their corresponding semivariance values grouped by distance. The distance is typically the Euclidean distance; here $d_{s_i}$ as defined in the common notation (see \autoref{sec:smoothingmethods}) is used. Semivariance describes the average squared difference between the values of points separated by a distance (or vector) $h$:
\[
\lambda_h=\frac{1}{2N(h)}\cdot\sum_{i=1}^{N(h)}{\left(z(x_i)-z(x_i+h)\right)}^2
\]

Because many pairs of points have unique distances, they are grouped into bins, similar to a histogram, and the average semivariance within each bin is used as the representative value. This results in a discrete function of distance that typically increases and then approaches a plateau. This reflects the empirical observation that nearby points tend to be more similar and the average difference stabilizes beyond a certain distance (see\ \cite{kriging:web01}).

The next step is to fit a mathematical model (usually denoted $\gamma(h)$) to the semivariogram. The models used follow the typical shape of a semivariogram, namely an increasing function that levels off at larger distances. The usual parameters are
\begin{itemize}
    \item $c_0$, the function baseline, also called the nugget,
    \item $c$, the value at which the function becomes approximately constant, and
    \item $a$, the distance at which this flattening occurs.
\end{itemize}

The simplest model is the linear model:
\[
\gamma(h)=
\begin{cases}
    c_0+c\left(\frac{h}{a}\right) & 0 < h \le a\\
    c_0+c & h > a\\
    c_0 & h=0
\end{cases}
\]
\begin{minipage}{\textwidth}
    The function $\gamma(h)$ is then used to compute the weights $\lambda_i$ in the smoothing formula (see \autoref{sec:smoothingmethods}). This is done by solving the Ordinary kriging system of equations
\end{minipage}
\[
\left(\begin{matrix}
    \gamma(x_1, x_1) & \dotsi & \gamma(x_1, x_N) & 1\\
    \vdots & \ddots & \vdots & \vdots\\
    \gamma(x_N, x_1) & \dotsi & \gamma(x_N, x_N) & 1\\
    1 & \dotsi & 1 & 0
\end{matrix}\right)
\left(\begin{matrix}
    \lambda_1\\
    \vdots\\
    \lambda_N\\
    \Psi
\end{matrix}\right)
=
\left(
\begin{matrix}
    \gamma(x_1,x_0)\\
    \vdots\\
    \gamma(x_N,x_0)\\
    1
\end{matrix}\right)
\]
where $\Psi$ is the Lagrange multiplier used to solve the system. The equation can also be written compactly as
\[
A\lambda=b
\]
solved for $\lambda$\ \cite{kriging:web02}
\[
\lambda=A^{-1}b
\]
which is the vector of all $\lambda_i$ for the smoothing formula.

\subsection{Shape Functions}\label{subsec:shape_functions}

Shape functions provide a way to interpolate a scalar or vector property within a small finite element by using the known values at its boundary nodes. In other words, they define how a value is distributed inside the element based on sampled data at the vertices. This section examines a Triangular Irregular Network (\emph{TIN}), following the descriptions in\ \cite{shape:book01} and\ \cite{shape:book02}.

A TIN is constructed by first generating a triangle mesh from a discrete set of data points using the Delaunay triangulation algorithm\ \cite{shape:triangles}. This ensures that no triangle contains any other data point and tends to avoid very slim triangles, which is good for numerical stability and output quality. After the triangulation has been formed, the interpolation starts with locating the specific triangle that contains the interpolation point. The interpolated value is then calculated as the weighted average of the values at the three vertices of the triangle.

\begin{figure}[H]
	\centering
    \includegraphics[scale=0.55]{images/triangle.png}
	\caption{A triangle within the TIN}\label{fig:TINtriangle}
\end{figure}

\autoref{fig:TINtriangle} illustrates a triangle formed by vertices $P_1$ through $P_3$, with the interpolation point within the triangle. The weight of each data point $P_i$ is determined by the area of the sub-triangle $A_i$ opposite of that vertex (or rather, the portion of the total area it covers).

The area of each sub-triangle is typically computed using the determinant formula
\[
A_t=\frac{1}{2}\det
\left(\begin{matrix}
    1 & x_1 & y_1\\
    1 & x_2 & y_2\\
    1 & x_3 & y_3
\end{matrix}\right)
\]
where $x_i$ and $y_i$ are the $x$- and $y$-coordinates, respectively, of the $i$-th point of that sub-triangle. To complete the interpolation, we define $\lambda_i$ as
\[
\lambda_i=\frac{A_i}{\sum_i^{N=3}A_i}
\]
so that
\[
\sum_{i=1}^{3}\lambda_i = 1
\]
for any point inside the triangle. Finally, the formula for the interpolated value $\hat{z}$ (see\ \autoref{sec:smoothingmethods}) at the query point $x_0$ adjusted to only use the three vertices of the triangle
\[
\hat{z}(x_0)=\sum_{i=1}^{N=3} \lambda_i\, z(x_i)
\]
\nolinebreak\cite{mdpi:01}

\subsection{Evaluation}\label{subsec:evaluation}
To compare the different smoothing methods, they will be tested on a real-world example dataset.

\subsubsection{Testing setup}
For this test, elevation data were chosen because they are easy to work with and accessible. For our testing purposes, the sector ``N46E15''\footnote{46° North, 15° East} of SRTM\footnote{Shuttle Radar Topography Mission} data was downloaded from \url{https://www.opendem.info}. It contains a 1° latitude by 1° longitude sector of the Earth.

The data come in the form of a \verb|.SHP|\footnote{Shapefile} file, which contains all contour lines and their heights. To convert them into a raster, the following Python script was used:
\begin{lstlisting}[language={[AS]Python},caption=Importing a raster from a Shapefile]
gdf = gpd.read_file("N46E015/N46E015.shp")
minx, miny, maxx, maxy = gdf.total_bounds
width = height = 500
transform = rasterio.transform.from_bounds(minx, miny, maxx, maxy, width, height)
shapes = ((geom, value) for geom, value in zip(gdf.geometry, gdf["elevation"]))

array = rasterio.features.rasterize(
    shapes=shapes,
    out_shape=(height, width),
    transform=transform,
    fill=np.nan,
    dtype="float32")

points = []
vals = []
for geom, val in zip(gdf.geometry, gdf["elevation"]):
    length = geom.length
    for i in np.linspace(0, length, 50):
        pt = shapely.ops.substring(geom, i, i)
        points.append((pt.x, pt.y))
        vals.append(val)

interpolated = scipy.interpolate.griddata(np.array(points), np.array(vals), tuple(np.meshgrid(
    np.linspace(minx, maxx, width),
    np.linspace(miny, maxy, height)
)), method="nearest")

mask = np.isnan(array)
array[mask] = np.flip(interpolated, 0)[mask]
\end{lstlisting}

\begin{itemize}
    \item \cdline{1} loads the data from the \verb|.shp| file.
    \item \cdlines{2--5} define constants and extract all shape objects from the file.
    \item \cdlines{7--13} use the \verb|rasterize| function from the \verb|rasterio| Python package to turn the shapes into a raster. This leaves holes in the array where there are no contour lines.
    \item \cdlines{15--22} set up the \verb|points| and \verb|vals| lists to be used in the interpolation by sampling points along each contour line.
    \item \cdlines{24--27} use the \verb|scipy| Python package to interpolate across a grid.
    \item \cdlines{29--30} replace all the holes in the array with interpolated values.
\end{itemize}

This produces the following height map:
\begin{figure}[H]
	\centering
    \includegraphics[scale=0.75]{images/N46E15.png}
	\caption{The extracted height map}\label{fig:heightmap}
\end{figure}

The different smoothing methods will now be tested by sampling 1000 data points from the height map, interpolating between them, and comparing the result with the original height map.

All implementations will be in the form of a function \verb|interpolate_<method>|, which takes the positions of the data points, the original height map, and additional parameters. All implementations will use the \verb|numpy| Python library to increase performance and provide the necessary operations.

\subsubsection{Nadaraya-Watson estimator}

\begin{minipage}{\textwidth}
The NW estimator is implemented using the following function:
\begin{lstlisting}[language={[AS]Python},caption=Interpolation with the NW estimator]
def interpolate_nwkr(points: np.ndarray, data: np.ndarray, kernel: Kernel, h: float):
    point_values = data[*points]
    out = np.empty(data.shape)
    for x in range(data.shape[0]):
        for y in range(data.shape[1]):
            dist_x = points[0].astype(np.float32) - x
            dist_y = points[1].astype(np.float32) - y
            density = kernel(np.hypot(dist_x, dist_y), h)
            out[x][y] = (density * point_values).sum() / density.sum()
    return out
\end{lstlisting}
\end{minipage}
\begin{itemize}
    \item \cdline{2} extracts the sample values from the original height map.
    \item \cdlines{4--5} loop over all interpolation points that make up the output grid.
    \item \cdlines{6--7} calculate the distance on the x and y axes from each data point to the interpolation point.
    \item \cdline{8} calculates the weight for each data point by passing the distance through the kernel. This is equivalent to $K_h\left(d_{s_i}\right)$.
    \item \cdline{9} calculates the interpolated value by multiplying each point weight by its value and dividing by the sum of all weights.
\end{itemize}

The function takes two additional parameters:
\begin{itemize}
    \item \verb|h|: the bandwidth of the kernel
    \item \verb|kernel|: the kernel function to use
\end{itemize}

As discussed in\ \autoref{subsec:nw_regression}, the NW estimator can use different kernel functions. This implementation provides the Gaussian and Epanechnikov kernels. Each kernel is implemented as a function that takes a numpy array of all data points and the bandwidth. This makes it possible to execute the function on many points at the same time.

\begin{minipage}{\textwidth}
The implementation of the Gaussian kernel:
\begin{lstlisting}[language={[AS]Python},caption=The Gaussian kernel]
def kernel(x: np.ndarray, h: float):
    u = x / h
    return 1 / np.sqrt(2 * np.pi) * np.exp(-np.square(u) / 2)
\end{lstlisting}
\end{minipage}

\begin{minipage}{\textwidth}
And the Epanechnikov kernel:
\begin{lstlisting}[language={[AS]Python},caption=The parabolic kernel]
def kernel(x: np.ndarray, h: float):
    u = x / h
    return np.maximum(3 / 4 * (1 - u * u), 0) / h
\end{lstlisting}
Note that values outside the bandwidth $|u|>1$ are handled by clipping any negative values to zero. This works because the parabolic kernel is negative outside of $[-1, 1]$.
\end{minipage}

\subsubsection{Inverse Distance Weighting}

IDW is implemented using the following function:
\begin{lstlisting}[language={[AS]Python},caption=Interpolation with IDW]
def interpolate_idw(points: np.ndarray, data: np.ndarray, u: float):
    point_values = data[*points]
    out = np.empty(data.shape)
    for x in range(data.shape[0]):
        for y in range(data.shape[1]):
            matches = (points[0] == x) & (points[1] == y)
            if matches.any():
                out[x][y] = point_values[matches].mean()
                continue

            dist_x = points[0].astype(np.float32) - x
            dist_y = points[1].astype(np.float32) - y
            dist_pow = np.power(np.hypot(dist_x, dist_y), -u)
            out[x][y] = (dist_pow * point_values).sum() / dist_pow.sum()

    return out
\end{lstlisting}
\begin{itemize}
    \item \cdline{2} extracts the sample values from the original height map.
    \item \cdlines{4--5} loop over all interpolation points that make up the output grid.
    \item \cdlines{6--9} check if any data points lie exactly on the interpolation point and use the mean of those that do as the interpolated value. This is necessary to avoid dividing by zero (see\ \autoref{subsec:idw}).
    \item \cdlines{11--12} calculate the distance on the x and y axes from each data point to the interpolation point.
    \item \cdline{13} calculates the weight for each data point by raising each distance to the power $-u$. This is equivalent to $K_h\left(d_{s_i}\right)$.
    \item \cdline{14} calculates the interpolated value by multiplying each point weight by its value and dividing by the sum of all weights.
\end{itemize}

The function takes the additional parameter \verb|u| that indicates how strongly the weight decreases as the distance increases.

\begin{minipage}{\textwidth}
Improved IDW can be achieved by introducing the parameter \verb|R| to the function and replacing \cdlines{13--14} with the following
\begin{lstlisting}[language={[AS]Python},caption=Addition for Improved IDW]
            dist = np.hypot(dist_x, dist_y)
            dist_pow = np.power(dist[dist <= R], -u)
            out[x][y] = (dist_pow * point_values[dist <= R]).sum() / dist_pow.sum()
\end{lstlisting}
\end{minipage}

\begin{itemize}
    \item \cdline{1} calculates the Euclidean distance from the x and y distances.
    \item \cdlines{2} calculates the distance raised to the power \verb|u| while selecting filtering for data points within the radius \verb|R|.
    \item \cdline{3} selects the values of the points within the radius \verb|R|.
\end{itemize}

\subsubsection{Distribution-based Distance Weighting}

DDW is implemented using the following function:
\begin{lstlisting}[language={[AS]Python},caption=Interpolation with DDW]
def interpolate_ddw(points: np.ndarray, data: np.ndarray, D: Distribution, R: float, k: float):
    point_values = data[*points]
    out = np.empty(data.shape)
    for x in range(data.shape[0]):
        for y in range(data.shape[1]):
            dist_x = points[0].astype(np.float32) - x
            dist_y = points[1].astype(np.float32) - y
            dist = np.hypot(dist_x, dist_y)
            values = point_values[dist <= R]
            mean = values.mean()
            sigma_R = R * values.std() / max(abs(np.min(values) - mean), abs(np.max(values) - mean))
            lambdas = D(dist[dist <= R], k * sigma_R)
            out[x][y] = (values * lambdas).sum() / lambdas.sum()
\end{lstlisting}

\begin{itemize}
    \item \cdline{2} extracts the sample values from the original height map.
    \item \cdlines{4--5} loop over all interpolation points that make up the output grid.
    \item \cdlines{6--7} calculate the distance on the x and y axes from each data point to the interpolation point.
    \item \cdline{8} calculates the Euclidean distance from the x and y distances.
    \item \cdline{9} selects the values of the points within the radius \verb|R|.
    \item \cdline{11} calculates $\sigma_R$ according to the formula (see \autoref{subsec:ddw}).
    \item \cdline{12} calculates all $\lambda_i$ by calling the distribution with the distances within the radius \verb|R|.
    \item \cdline{13} calculates the interpolated value by multiplying each point weight by its value and dividing by the sum of all weights.
\end{itemize}

The function takes three additional parameters:
\begin{itemize}
    \item \verb|D|: the distribution to use
    \item \verb|R|: the radius of the circle in which to consider the data points
    \item \verb|k|: the custom parameter by which to scale $\sigma_R$
\end{itemize}

The three distributions discussed in this thesis (see \autoref{subsec:ddw}) are implemented as Python functions that take a numpy array of distances and their corresponding spread value. Because $\mu_R$ is 0 for our purpose, $\mu$ was taken as 0 and removed from all distributions.

\begin{minipage}{\textwidth}
The implementation of the normal distribution:
\begin{lstlisting}[language={[AS]Python},caption=The normal distribution for DDW]
def dist_normal(x: np.ndarray, sigma: float):
    return 1 / (np.sqrt(2 * np.pi) * sigma) * np.exp(-np.square(x) / (2 * sigma * sigma))
\end{lstlisting}
\end{minipage}

\begin{minipage}{\textwidth}
The implementation of the Cauchy distribution:
\begin{lstlisting}[language={[AS]Python},caption=The Cauchy distribution for DDW]
def dist_cauchy(x: np.ndarray, gamma: float):
    return gamma / (np.pi * (np.square(x) + gamma * gamma))
\end{lstlisting}
\end{minipage}

\begin{minipage}{\textwidth}
The implementation of the Laplace distribution:
\begin{lstlisting}[language={[AS]Python},caption=The Laplace distribution for DDW]
def dist_laplace(x: np.ndarray, b: float):
    return 1 / (2 * b) * np.exp(-np.abs(x) / b)
\end{lstlisting}
\end{minipage}

\subsubsection{Triangular Irregular Network}

This example uses the \verb|matplotlib| Python package for interpolation:
\begin{lstlisting}[language={[AS]Python},caption=Interpolation with a TIN]
def interpolate_tin(points: np.ndarray, data: np.ndarray):
    point_values = data[*points]
    triangulation = matplotlib.tri.Triangulation(*points)
    interpolator = matplotlib.tri.LinearTriInterpolator(triangulation, point_values)
    interp = interpolator(*np.meshgrid(np.arange(data.shape[0]), np.arange(data.shape[1])))
    return np.array(interp.filled(np.nan))
\end{lstlisting}

\begin{itemize}
    \item \cdline{2} extracts the sample values from the original height map.
    \item \cdline{3} creates a triangulation using the Delaunay triangulation algorithm (see \autoref{subsec:shape_functions}).
    \item \cdline{4} creates the interpolator object using the data points.
    \item \cdline{5} runs the interpolator on a grid.
\end{itemize}

\subsubsection{Kriging}
\enlargethispage{1pt}
This example uses the \verb|pykrige| Python package for kriging:
\begin{lstlisting}[language={[AS]Python},caption=Interpolation with Kriging]
def interpolate_kriging(points: np.ndarray, data: np.ndarray, bins: float, model: str):
    kriging = pykrige.ok.OrdinaryKriging(x=points[0],
                                         y=points[1],
                                         z=data[*points],
                                         nlags=bins,
                                         variogram_function=model)
    return kriging.execute("grid", np.arange(data.shape[0]), np.arange(data.shape[1]))
\end{lstlisting}

The function takes two additional parameters:
\begin{itemize}
    \item \verb|bins|: the number of lag bins used when computing the semivariogram
    \item \verb|model|: the model used to fit the semivariogram
\end{itemize}

\subsubsection{Comparison}

To provide an objective measure of the quality of each smoothing method, the RMSE\footnote{Root Mean Squared Error} is used. It is calculated as follows:
\[
RMSE=\sqrt{\frac{1}{n}\sum_{i=1}^N{\left(z(x_i)-\hat{z}(x_i)\right)}^2}
\]
\nolinebreak\cite{shape:book01} where $\hat{z}(i)$ is the interpolated value and $z(i)$ is the actual value at a data point $i$.

The RMSE is calculated using the following Python function:
\begin{lstlisting}[language={[AS]Python},caption=Calculate RMSE]
def rmse(a: np.ndarray, b: np.ndarray):
    return np.sqrt(np.square(a - b).mean())
\end{lstlisting}

\begin{table}[H]
    \centering
    \caption{Comparison of Smoothing Methods}
    \begin{tabular}{| p{18em} | c | c | >{\centering\arraybackslash}p{9em} |} % chktex 44
        \hline % chktex 44
        Method & Runtime & RMSE & Notes \\
        \hline % chktex 44
        NW estimator using the Gaussian kernel, h~=~10 & 4.8s & 91.77 & Smallest h that doesn't leave holes\\
        NW estimator using the Parabolic kernel, h~=~35 & 4.2s & 102.71 & Smallest h that doesn't leave holes\\
        IDW with u = 3.75 & 5.2s & 89.46 & \\
        Improved IDW with u = 3.5, R = 40 & 4.8s & 88.71 & No spatial data structure was used \\
        DDW with Normal distribution, R = 50, k = 0.5 & 9s & 95.06 & \\
        DDW with Cauchy distribution, R = 50, k = 0.35 & 9s & 99.2 & \\
        DDW with Laplace distribution, R = 50, k = 0.35 & 9s & 93.60 & \\
        Kriging with 64 bins and spherical model & 3m & 96.52 & \\
        TIN & 50ms & 93.03 & Data points included corners \\
        \hline % chktex 44
    \end{tabular}
\end{table}

Note:
\begin{itemize}
    \item Air-quality data and elevation data differ dramatically. However, elevation data were easier to obtain and are considered sufficient for this comparison.
    \item This comparison is not a complete assessment of each interpolation method. Its purpose is to provide an overview of the methods and their performance and visual results. For a more sophisticated comparison, see\ \cite{shape:book01}.
    \item The positions of the data points were chosen randomly. A proper comparison should use points that are uniformly distributed to a certain degree (see\ \cite{mdpi:01}). However, our data will be unevenly distributed, and the handling of uneven point density is a major point of this thesis discussed in a later chapter.
    \item The parameters of the methods were found by experimentation. This is because most parameters do not have a way of obtaining an optimal value (see $R$ in \autoref{subsec:idw}). One way would be to use function optimization (see\ \cite{other:scipy_optimization}) to minimize RMSE, but that is outside the scope of this thesis.
\end{itemize}

\begin{figure}[ht]
\centering
\captionsetup{justification=centering}

\begin{minipage}[t]{0.32\textwidth}
  \centering
  \includegraphics[width=\linewidth]{images/interpolations/nwkr_gauss_10.png}
    \caption{NW estimator using the Gaussian kernel}
\end{minipage}
\hfill
\begin{minipage}[t]{0.32\textwidth}
  \centering
  \includegraphics[width=\linewidth]{images/interpolations/nwkr_parab_35.png}
    \caption{NW estimator using the Parabolic kernel}
\end{minipage}
\hfill
\begin{minipage}[t]{0.32\textwidth}
  \centering
  \includegraphics[width=\linewidth]{images/interpolations/kriging_64_spherical.png}
  \caption{Kriging}
\end{minipage}

\end{figure}

\begin{figure}[H]
\centering
\captionsetup{justification=centering}

\begin{minipage}[t]{0.32\textwidth}
  \centering
  \includegraphics[width=\linewidth]{images/interpolations/idw_3.75.png}
  \caption{IDW}
\end{minipage}
\hfill
\begin{minipage}[t]{0.32\textwidth}
  \centering
  \includegraphics[width=\linewidth]{images/interpolations/idw_imp_3.5_40.png}
  \caption{Improved IDW}
\end{minipage}
\hfill
\begin{minipage}[t]{0.32\textwidth}
  \centering
  \includegraphics[width=\linewidth]{images/interpolations/tin.png}
  \caption{TIN}
\end{minipage}

\end{figure}

\begin{figure}[H]
\centering
\captionsetup{justification=centering}

\begin{minipage}[t]{0.32\textwidth}
  \centering
  \includegraphics[width=\linewidth]{images/interpolations/ddw_normal_50_.5.png}
  \caption{DDW with Normal distribution}
\end{minipage}
\hfill
\begin{minipage}[t]{0.32\textwidth}
  \centering
  \includegraphics[width=\linewidth]{images/interpolations/ddw_cauchy_50_.35.png}
  \caption{DDW with Cauchy distribution}
\end{minipage}
\hfill
\begin{minipage}[t]{0.32\textwidth}
  \centering
  \includegraphics[width=\linewidth]{images/interpolations/ddw_laplace_50_.35.png}
  \caption{DDW with Laplace distribution}
\end{minipage}
\end{figure}

\section{Reference Implementation using Python}

To implement a value distribution map, the following procedure is used:
\begin{itemize}
    \item Create an empty overlay to display on the map.
    \item For every pixel, use a smoothing method to get the interpolated value at that pixel and use a color map to obtain the color for that pixel.
    \item For every pixel, use a kernel density estimation to get the opacity of the overlay.
\end{itemize}

The following implementation uses an NW estimator as the smoothing method, with the same kernel and bandwidth as the KDE\@.

\begin{lstlisting}[language={[AS]Python},caption=Reference implementation in Python]
BANDWIDTH = 20

def kernel(x):
    return np.exp(-1 * x * x)

def smoothstep(x):
    return x * x * (3 - 2 * x)

def kve(points_x, points_y, values, x, y):
    distance = np.hypot(points_x - x, points_y - y)
    density = np.sum(kernel(distance / BANDWIDTH))
    value = np.sum(kernel(distance / BANDWIDTH) * values)

    return (density / (BANDWIDTH * len(points_x)), value / density if density > 0 else 0)

def cmap(value):
    return (193 + value * 62, 88 + value * 167, value * 255)

def cmult(color, mult):
    return (int(color[0] * mult), int(color[1] * mult), int(color[2] * mult))

def draw_kve(pxarray: np.ndarray, X: np.ndarray, Y: np.ndarray, VALUES: np.ndarray):
    densities = np.zeros((width, height), dtype=np.float32)
    values = np.zeros((width, height), dtype=np.float32)

    for x in range(width):
        for y in range(height):
            density, weight = kve(X, Y, VALUES, x, y)
            densities[x, y] = density
            values[x, y] = weight

    max_density = np.max(densities)

    for x in range(width):
        for y in range(height):
            density = densities[x, y]
            value = values[x, y]
            density = smoothstep(density / max_density) if max_density > 0 else 0
            color = cmap(value / 255)
            pxarray[x, y] = cmult(color, density)
\end{lstlisting}
\begin{itemize}
    \item \cdlines{1--7} define the bandwidth, the kernel function (here the Gaussian kernel is used), and the \verb|smoothstep| function\cite{quilez:smoothstep}, which smoothly interpolates values from 0 to 1.
    \item \cdlines{9--14} define the \verb|kve| function, which calculates both the KDE and the NW estimator at a point $(x, y)$.
    \item \cdlines{16--17} define the \verb|cmap| function, which maps a value from 0 to 1 to a color, here from brown to white.
    \item \cdlines{19--20} define the \verb|cmult| function, a utility for multiplying all components of an RGB color by a scalar.
    \item \cdlines{22--40} define the function \verb|draw_kve|, which draws the results onto a pixel array.
    \item \cdlines{26--30} loop over all pixels, calculating all densities and interpolated values.
    \item \cdline{38} normalizes the density by dividing by the maximum density and applies the smoothstep function.
    \item \cdlines{39--40} apply the color map and multiplies by the normalized density. Note that in this example, a black background is used instead of transparency.
\end{itemize}

The implementation produces the following output:
\begin{figure}
    \centering
    \includegraphics[scale=0.5]{images/Interp_Python.png}
    \caption{An example output of the Python implementation}\label{fig:interp_python}
\end{figure}

\subsection{Flaws}\label{subsec:flaws}
This implementation has the following flaws:
\begin{itemize}
    \item The data values will likely not range from 0 to 255
    \item Because density is normalized on the basis of the maximum, an area with high point density reduces the visibility of all other data points and could in extreme cases fade out other points completely.
    \item The density is also divided by the total number of points. This aligns with the definition of the KDE\@; however, the density of a local area can be reduced if data points are introduced at a distant location.
\end{itemize}

These flaws will be addressed in the final implementation at the end of this thesis.